% Start preamble
\documentclass[12pt,a4paper]{article}
\usepackage{geometry}
 \geometry{
 a4paper,
 total={170mm,257mm},
 left=20mm,
 top=20mm,
 }
\usepackage{currfile}
\usepackage[utf8]{inputenc}
\usepackage[T1]{fontenc}
\usepackage[pdftex]{graphicx}
\graphicspath{{./}}
\usepackage{enumitem}
\usepackage{pdfpages}
\usepackage{hyperref}
\usepackage{tikz}
\usepackage{attachfile}
\usepackage{epstopdf}
\usepackage{array}
\usepackage{multirow}
\usepackage{multicol}
\usepackage{float}
%\usepackage[table]{xcolor,colorbl}
\setlength{\textwidth}{16cm}
\setlength{\oddsidemargin}{-0.5cm}
\setlength{\evensidemargin}{-0.5cm}
%\setlenght{\headsep}{0cm}
\setlength\parindent{0pt}
%\setlength{\extrarowheight}{3pt}
\usepackage{listings}
%\usepackage{xcolor}

\input{arduinoLanguage.tex}
%%%%%% Counting oppgaves %%%%%%
 \newcount\questnum \questnum=0
 \def\oppgave{

		}

% End preamble

\begin{document}
\title{TIF Prosessforståelse 01}
\author{Faglærer: Fred-Olav Mosdal 90507684\\}
\maketitle
\oppgave{}%1
\vskip 2.5pt 

Hva er en prosess? Forklar med et praktisk eksempel. 


\vskip 5pt 
\begin{tikzpicture}
	\draw[step=0.5cm,gray!20,very thin]  grid (17,6) ;
\end{tikzpicture}
\vskip 2.5pt 
\oppgave{}%1
\vskip 2.5pt 

Hva er en kjemisk prosess?


\vskip 5pt 
\begin{tikzpicture}
	\draw[step=0.5cm,gray!20,very thin]  grid (17,6) ;
\end{tikzpicture}
\vskip 2.5pt 
\oppgave{}%1
\vskip 2.5pt 

Hva mener vi med et «pådrag» i en prosess?


\vskip 5pt 
\begin{tikzpicture}
	\draw[step=0.5cm,gray!20,very thin]  grid (17,6) ;
\end{tikzpicture}
\vskip 2.5pt 
\oppgave{}%1
\vskip 2.5pt 

Hva mener vi med en «utgang» i en prosess?


\vskip 5pt 
\begin{tikzpicture}
	\draw[step=0.5cm,gray!20,very thin]  grid (17,6) ;
\end{tikzpicture}
\vskip 2.5pt 
\oppgave{}%1
\vskip 2.5pt 

Hva menes med «samvirkninger» i en prosess?


\vskip 5pt 
\begin{tikzpicture}
	\draw[step=0.5cm,gray!20,very thin]  grid (17,6) ;
\end{tikzpicture}
\vskip 2.5pt 
\oppgave{}%1
\vskip 2.5pt 

Hva er «god» og «dårlig» kvalitet på et produkt?


\vskip 5pt 
\begin{tikzpicture}
	\draw[step=0.5cm,gray!20,very thin]  grid (17,6) ;
\end{tikzpicture}
\vskip 2.5pt 
\oppgave{}%1
\vskip 2.5pt 

Hva betyr «systematisk», og hvorfor må operasjonene i en prosess utføres på en systematisk måte?


\vskip 5pt 
\begin{tikzpicture}
	\draw[step=0.5cm,gray!20,very thin]  grid (17,6) ;
\end{tikzpicture}
\vskip 2.5pt 
\oppgave{}%1
\vskip 2.5pt 

Å bake brød er en kjemisk prosess. Et enkelt brød kan du lage ved å blande ingrediensene mel, gjær og vann. Deretter må deigen heve, brødet må bakes ut, og til slutt må det stekes.

a Tegn et enkelt flytskjema som viser enhetsoperasjonene og sammenhengen mellom disse.
\vskip 0.5 cm 
b Mengden mel og steketemperatur er noen avpådragene i denne prosessen. Kan du finne flere?
\vskip 0.5 cm 
c Smaken på brødet er en av utgangene i denne prosessen. Kan du finne flere?
\vskip 0.5 cm 
d Hva slags prosessforstyrrelser tror du kan inntreffe?
\vskip 0.5 cm 

\vskip 5pt 
\begin{tikzpicture}
	\draw[step=0.5cm,gray!20,very thin]  grid (17,16) ;
\end{tikzpicture}
\vskip 2.5pt 
\end{document}

